\needspace{8\baselineskip}
\section{模型假设}

下列假设只补充题面未写明、但求解时必须固定的条款；题面已经给出的几何、计时与接口规则不再重复列为假设。假设的作用是把最坏界、外包络与在线信息集写死，避免后文在问题2用真值排点、在问题4把无信号读成距离圆盘。

\begin{enumerate}
\item \textbf{平面与直线。} 忽略测向机与干扰源的高度差，全部对象落在同一平面；机器狗在两次停测之间沿直线匀速运动。题面已声明高度差可忽略，直线是接口允许的最短可行路径，不引入转弯半径或地形障碍。坐标分量绝对值不得超过 $2\times 10^6\,\mathrm{m}$，越界点直接丢弃。圆域 $\Omega$ 只约束源位置，不禁止机器狗出圆。

\item \textbf{误差的最坏界，而非随机分布。} 示向度误差落在 $[-1^\circ,1^\circ]$ 闭区间，同一地点同一频道的误差固定。本文不指定误差在区间内的概率分布，定位区域按最坏界取闭楔形。接口将示向度保留两位小数，实现中用 $1.005^\circ$ 作为数值包络以吸收量化，物理误差常数仍取题设的 $1^\circ$，该包络不进入参数表，也不进入问题1按题面作答的半宽。

\item \textbf{接收半径未知且异质。} 各源 $r_i$ 独立落在 $[1000,1500]\,\mathrm{m}$，接口不返回其值。凡声称“一定收得到信号”的几何结论，一律用最紧下界 $1000\,\mathrm{m}$；凡声称“可能收不到”的外包络，一律用 $1500\,\mathrm{m}$。不得用一次 $\mathtt{no\_signal}$ 反推精确距离。问题4的无信号增加扇外分支，不得挖掉距离圆盘。

\item \textbf{覆盖与示向的方向约定。} 射频覆盖用源指向检测点的方向判定；示向度是检测点指向源的方位。定向源覆盖为朝向两侧各 $90^\circ$（含）。光学清除只检查欧氏距离，与射频扇无关。该约定与题面附录一致，写在此处是为了后文半平面证明不再每次重述。$5\,\mathrm{m}$ 与 $20\,\mathrm{m}$ 均为闭边界。

\item \textbf{在线决策不使用真值。} 机器狗只能使用自身指令的反馈。演练结束返回的 $N$、离线实验里的源坐标与朝向，一律不得写入在线策略。离线几何模型仅用于配对比较与回归，不替代模拟器正式测试。正式测试永不返回 $N$，不能把 $K=10$ 写成停机。启发式预算不得升级为永久放弃。
\end{enumerate}

上述五条把后文反复用到的边界钉死。平面假设使闭楔形成为二维半平面交，而不是三维圆锥交；若将来测向机提供俯仰角，问题1的对象必须改写，不能把平面直径直接当作空间直径。直线假设与接口一致：两次停测之间没有转弯半径，虚拟时间按欧氏距离除以 $5\,\mathrm{m/s}$ 计入，因此寻路层用插入阈值约束的是额外路程，而不是曲率。

最坏界假设决定了问题1不用协方差椭圆。同一点同一频道的误差固定，重复测量不能把楔形半宽从 $1^\circ$ 削到 $0.5^\circ$。数值包络 $1.005^\circ$ 只吸收两位小数量化，不得写成“题面其实允许一点二度”。接收半径异质假设决定了覆盖命题一律用 $1000\,\mathrm{m}$ 下界，外包络一律用 $1500\,\mathrm{m}$ 上界；一次无信号既不能反推精确 $r_i$，在问题4里也不能挖掉距离圆盘。

覆盖与示向相差 $180^\circ$，写进假设是为了半平面法向量的指向不再每次重述。
光学半径与射频扇解耦，则扇外 $d\le 20\,\mathrm{m}$ 仍可清除，问题4的光学方格才有独立地位。
在线不使用真值，则演练弹窗里的 $N$ 只能事后核对 $K/N$，离线种子里的朝向只能出现在存在性证明，不能出现在 $\mathtt{/measure}$ 的选点函数里。
启发式预算一旦写成永久放弃，字典序目标就被破坏，问题3的 ABORT 对照已经给出反例：全清率降至 $1/10$，平均 $T/K=382\,\mathrm{s}$ 并不能补偿漏清。
五条假设因此不是修辞，而是后文证书能否成立的前提：平面给出二维半平面，最坏界给出闭楔形，异质半径给出 $1000$ 与 $1500$ 的分工，方向约定给出半平面法向量，在线信息集禁止把真值写回策略。
